\chapter{Investigação}
\label{chap:investigacao}
\label{chap:resultados-discussao}

Este capítulo descreve a metodologia adotada para a execução de experimentos e apresenta os resultados
obtidos. 

O objetivo em nenhum momento é demonstrar
vantagem quântica, mas sim avaliar se o sistema produz soluções válidas, como se comporta
sob diferentes configurações e onde se encontram os seus limites.

A execução dos experimentos foi assistida por um agente de inteligência artificial 
com acesso ao repositório, que implementou os cenários definidos, correu as 
simulações e gerou os relatórios de resultados. Todas as conclusões foram 
posteriormente revistas e validadas pelo autor. No total foram realizadas 1164
execuções principais, distribuídas por cinco fases experimentais.
\section{Arquitetura dos Experimentos}
\label{sec:arquitetura-experimentos}

Para avaliar o sistema em condições distintas, foi criada uma camada experimental em
torno do \textit{SchedulingEngine}. Esta camada separa a definição dos cenários da sua
execução, o que permite variar sistematicamente a dimensão do problema, o número de
processos, o limite de qubits, os parâmetros do QAOA e a estratégia de resolução sem
alterar o código principal do sistema.

A arquitetura experimental organiza-se em quatro etapas: configuração do cenário,
orquestração da execução, resolução pela pipeline selecionada e persistência dos
resultados. A Figura~\ref{fig:fluxo-experimentos} apresenta este fluxo. Os executores de
cenários constituem a camada adicional sobre a arquitetura descrita no capítulo anterior;
o restante processamento continua a ser realizado pelo \textit{SchedulingEngine}.

\begin{figure}[H]
\centering
\begin{tikzpicture}[
    node distance=1.4cm,
    box/.style={rectangle, draw, rounded corners, align=center, minimum width=3.7cm, minimum height=0.9cm},
    arrow/.style={->, thick}
]

\node[box] (toml) {Ficheiro TOML\\do cenário};
\node[box, below of=toml] (runner) {Scenario Runner\\ou Sweep Runner};
\node[box, below of=runner] (engine) {Scheduling Engine};
\node[box, below left=1.2cm and 0.8cm of engine] (default) {Pipeline direta};
\node[box, below right=1.2cm and 0.8cm of engine] (iterative) {Pipeline iterativa\\com sub-QUBOs};
\node[box, below=2.4cm of engine] (solver) {Solver QAOA\\e validação};
\node[box, below of=solver] (results) {Resultados\\JSON/JSONL, TXT e logs};

\draw[arrow] (toml) -- (runner);
\draw[arrow] (runner) -- (engine);
\draw[arrow] (engine) -- (default);
\draw[arrow] (engine) -- (iterative);
\draw[arrow] (default) -- (solver);
\draw[arrow] (iterative) -- (solver);
\draw[arrow] (solver) -- (results);

\end{tikzpicture}
\caption{Arquitetura utilizada na execução dos experimentos da investigação.}
\label{fig:fluxo-experimentos}
\end{figure}

\subsection{Definição e Execução dos Cenários}

Cada experimento é descrito num ficheiro de configuração em formato TOML. As secções
do ficheiro definem o \textit{workload}, a formulação QUBO, os parâmetros do QAOA, a
decomposição, a validação e as opções de execução. Esta separação permite repetir uma
experiência ou alterar uma única variável experimental sem modificar o código do
sistema.

O \textit{Scenario Runner} carrega a configuração, converte os seus valores nos objetos
usados pelo sistema e executa o número de repetições especificado. Nos estudos
paramétricos, o \textit{Sweep Runner} expande os eixos e casos definidos no ficheiro TOML,
gera as combinações correspondentes e entrega cada variante ao mesmo mecanismo de
execução. Deste modo, tanto uma execução isolada como uma varredura de parâmetros seguem
o mesmo percurso e produzem registos com a mesma estrutura.

Para cada variante, o \textit{SchedulingEngine} seleciona a pipeline direta ou a pipeline
iterativa. A primeira resolve a instância QUBO completa; a segunda decompõe o problema em
sub-QUBOs quando o limite configurado de qubits é ultrapassado. A solução produzida pelo
solver é depois validada e resumida num conjunto comum de métricas.

Um exemplo simplificado de um cenário experimental é apresentado de seguida:

\begin{lstlisting}[caption={Exemplo simplificado de um cenário experimental em TOML.}]
scenario_version = 1
id = "T3.5"
name = "Three-core iterative decomposition"

[workload]
mode = "preset"
num_processes = 5
num_cores = 3
weights = [0.36, 0.27, 0.22, 0.18, 0.12]

[qubo]
penalty = 5.0

[qaoa]
layers = 1
steps = 50
learning_rate = 0.05
top_k = 20
mixer_type = "xy"

[decomposition]
qubit_max = 9
num_cores = 3
sorting_strategy = "WEIGHT_DESCENDING"

[execution]
pipeline = "iterative"
solvers = ["qaoa"]
repeats = 1
\end{lstlisting}

\subsection{Recolha e Armazenamento dos Resultados}

Após cada execução, o sistema cria uma pasta própria para o cenário e guarda uma
cópia da configuração utilizada. Cada repetição origina um ficheiro JSON com a
configuração resolvida e as métricas, um ficheiro TXT legível e um log da saída produzida.

Em cada execução são registados, entre outros, os seguintes dados:

\begin{itemize}
    \item configuração usada no experimento;
    \item número de processos, núcleos e variáveis QUBO;
    \item energia da solução obtida;
    \item viabilidade da solução;
    \item tempo de resolução;
    \item atribuição final de processos a cores;
    \item resultado da validação;
    \item desequilíbrio final de carga entre os núcleos;
\end{itemize}


\clearpage
\section{Resultados da Investigação}

\subsection{Validação Inicial da Pipeline}

Foram inicialmente executadas instâncias pequenas, com dois a quatro processos e dois
cores. Estas instâncias produzem entre 4 e 8 variáveis QUBO, o que permite calcular o
ótimo global por \textit{brute force} e comparar diretamente a solução obtida pelo QAOA.

\begin{table}[H]
\centering
\begin{tabular}{|c|c|c|c|c|c|}
\hline
\textbf{Processos} & \textbf{Cores} & \textbf{Variáveis} & \textbf{Viável} & \textbf{Ótimo} & \textbf{Pipeline} \\
\hline
2 & 2 & 4 & Sim & Sim & Direta \\
3 & 2 & 6 & Sim & Sim & Direta \\
4 & 2 & 8 & Sim & Sim & Direta \\
\hline
\end{tabular}
\caption{Resultados dos testes iniciais em instâncias pequenas.}
\label{tab:resultados-sanity}
\end{table}

Em todos os casos a solução produzida foi viável e coincidiu com o ótimo obtido por
\textit{brute force}. Isto confirma que, para instâncias pequenas, o sistema executa corretamente
o ciclo completo.

Numa instância com quatro processos e dois cores, foi também avaliado o efeito da
profundidade $p$ do circuito QAOA, ou seja, o número de vezes que os operadores de fase
e de mistura são aplicados. Foram testados valores $p \in \{1, \ldots, 7\}$ com cinco
repetições cada. Em todos os casos $p=1$ foi suficiente para atingir o ótimo, com o
tempo a crescer de 1277\,ms para 5111\,ms sem qualquer melhoria de qualidade; esta
instância era demasiado simples para discriminar o efeito da profundidade, cujo impacto
real só se torna visível em instâncias maiores.

\subsubsection{Validação da Decomposição}

Para avaliar o comportamento da pipeline à medida que o número de processos cresce, foram
executadas instâncias com $N \in \{2, 3, 4, 5, 6, 8, 10\}$ processos e dois cores, com
\texttt{qubit\_max}=6. Quando o número de variáveis QUBO excede este limite, o problema
é automaticamente dividido em sub-QUBOs resolvidos pela pipeline
iterativa.

\begin{table}[H]
\centering
\resizebox{\textwidth}{!}{%
\begin{tabular}{|c|c|c|c|c|c|c|}
\hline
\textbf{$N$} & \textbf{Variáveis} & \textbf{Pipeline} & \textbf{Sub-QUBOs} & \textbf{Ótimo} & \textbf{Load imbalance} & \textbf{Tempo médio (ms)} \\
\hline
2  & 4  & Direta    & 0 & 100\% & 0.000 & 896.1  \\
3  & 6  & Direta    & 0 & 100\% & 0.333 & 1279.6 \\
4  & 8  & Iterativa & 2 & 100\% & 0.000 & 2401.3 \\
5  & 10 & Iterativa & 2 & 100\% & 0.200 & 2869.1 \\
6  & 12 & Iterativa & 2 & 100\% & 0.000 & 3279.0 \\
8  & 16 & Iterativa & 3 & 100\% & 0.000 & 4202.9 \\
10 & 20 & Iterativa & 4 & 100\% & 0.000 & 5755.9 \\
\hline
\end{tabular}
}
\caption{Escalabilidade da pipeline com o número de processos.}
\label{tab:resultados-d1}
\end{table}

A decomposição manteve a viabilidade e atingiu o
ótimo global em todas as execuções até $N=10$. O tempo cresceu com o número de
sub-QUBOs, de 896\,ms em $N=2$ para 5756\,ms em $N=10$.

Os desequilíbrios de carga em $N=3$ e $N=5$ não resultam de falha da pipeline: com pesos
uniformes e dois cores, é impossível distribuir um número ímpar de processos sem algum
desequilíbrio.

\subsubsection{Caracterização Inicial do QAOA}
\label{subsec:caracterizacao-inicial-qaoa}
Foi utilizada uma instância fixa com oito processos e dois núcleos
($N=8$, $K=2$), correspondendo a 16 variáveis QUBO. Os pesos dos processos foram
$$
\mathbf{w} = [0.30,\,0.27,\,0.22,\,0.18,\,0.15,\,0.11,\,0.08,\,0.05],
$$

O parâmetro \texttt{qubit\_max} define o número máximo de variáveis que o QAOA pode 
resolver de uma vez; quando o problema excede esse limite, é dividido em sub-QUBOs 
resolvidos sequencialmente. 

A profundidade $p$ controla o número de camadas alternadas 
de operadores de fase e de mistura no circuito --- maior profundidade aumenta a 
expressividade do circuito mas também o custo de simulação. 

Os passos referem-se ao 
número de iterações do otimizador clássico que ajusta os parâmetros $\gamma$ e $\beta$ 
do circuito.

Cada resolução usa um
circuito de profundidade $p=2$ e 100 passos do otimizador clássico.

\begin{table}[H]
\centering
\begin{tabular}{|c|c|c|c|c|}
\hline
\textbf{\texttt{qubit\_max}} & \textbf{Pipeline} & \textbf{Sub-QUBOs} & \textbf{Gap médio} & \textbf{Tempo médio (ms)} \\
\hline
4  & Iterativa & 4 & 0.000000 & 4550.3 \\
6  & Iterativa & 3 & 0.000000 & 4546.8 \\
8  & Iterativa & 2 & 0.000000 & 4777.0 \\
10 & Iterativa & 2 & 0.000000 & 4792.5 \\
12 & Iterativa & 2 & 0.000000 & 5648.0 \\
16 & Direta    & 0 & 0.000591 & 6911.1 \\
\hline
\end{tabular}
\caption{Comparação entre execução iterativa e direta para diferentes limites de qubits ($N=8$, $K=2$).}
\label{tab:resultados-qubitmax}
\end{table}

Nos casos em que o limite de qubits forçou a decomposição, a pipeline iterativa atingiu o ótimo global em todas as execuções. Com \texttt{qubit\_max}=16,
a pipeline resolveu a instância completa sem decomposição, mas não atingiu o ótimo com
o orçamento de $p=2$ e 100 passos, uma vez que o solver passou a otimizar um espaço
de estados de maior dimensão com o mesmo orçamento algorítmico.

Este resultado sugere que a configuração $p=2$ e 100 iterações não foi suficiente
para explorar adequadamente o espaço de estados da instância. Para testar
esta interpretação, manteve-se \texttt{qubit\_max}=16 e variaram-se a profundidade
e o número de iterações do otimizador.

\begin{table}[H]
\centering
\begin{tabular}{|c|c|c|c|c|}
\hline
\textbf{Profundidade $p$} &
\textbf{Iterações} &
\textbf{Gap} &
\textbf{Ótimo} &
\textbf{Tempo (ms)} \\
\hline
2 & 100 & 0.000591 & Não & 8087.1  \\
2 & 200 & 0.000591 & Não & 14141.7 \\
2 & 500 & 0.001251 & Não & 35510.7 \\
\hline
3 & 100 & 0.000000 & Sim & 10577.5 \\
3 & 200 & 0.000000 & Sim & 20738.6 \\
3 & 500 & 0.000000 & Sim & 51943.0 \\
\hline
4 & 100 & 0.000000 & Sim & 13701.7 \\
4 & 200 & 0.000000 & Sim & 27484.1 \\
4 & 500 & 0.000000 & Sim & 68555.0 \\
\hline
\end{tabular}
\caption{Impacto da profundidade e do número de iterações na resolução direta
com \texttt{qubit\_max}=16 ($N=8$, $K=2$).}
\label{tab:resultados-profundidade-passos-direta}
\end{table}
O aumento da profundidade de $p=2$ para $p=3$ permitiu recuperar o ótimo global
com apenas 100 iterações. Em contraste, aumentar exclusivamente o número de
iterações, mantendo $p=2$, não melhorou a solução. Aliás, com 500 iterações, o gap foi
inclusivamente superior ao observado com 100 e 200 iterações.

Estes resultados indicam que, nesta instância, a profundidade do circuito teve
maior influência na qualidade da solução do que o número de iterações do
otimizador. A resolução direta conseguiu, portanto, atingir o ótimo global, mas
exigiu uma profundidade superior e maior tempo computacional.

No entanto, a pipeline iterativa atingiu a mesma qualidade em menos tempo nas configurações
testadas. Este resultado demonstra uma vantagem computacional para esta instância e mostra que o
aumento do limite de qubits deve ser acompanhado por uma parametrização adequada do QAOA.

\paragraph{Caracterização do Otimizador Clássico.}
Foi também avaliado o efeito do número de iterações do otimizador clássico e da 
inicialização de $\gamma/\beta$. Testaram-se valores entre 50 e 500 iterações e um 
grid de 25 pares de inicialização; em todos os casos a solução foi viável e ótima. 
Aumentar de 50 para 500 iterações multiplicou o tempo por um fator de aproximadamente 
10 sem qualquer melhoria de qualidade, o que reforça que estas instâncias eram 
demasiado simples para revelar o limite do otimizador.

\clearpage
Os resultados da Secção~\ref{subsec:caracterizacao-inicial-qaoa} mostram que a
qualidade depende não só do orçamento do circuito, mas também da forma como a pipeline
gere a informação entre sub-QUBOs. As duas secções seguintes caracterizam os parâmetros
que controlam essa gestão.

\section{Escolha do Mixer}

O mixer X aplica rotações independentes a cada qubit e não preserva a restrição 
de atribuição única durante a evolução do circuito, dependendo da penalidade $P$ para 
desencorajar estados inválidos. O mixer XY atua apenas entre qubits do mesmo 
processo, preservando essa restrição estruturalmente e eliminando essa dependência.

O valor de $P_\text{safe}$ é calculado pela heurística
\[
P_\text{safe} \approx 2 \cdot w_{\max} \cdot \frac{W_\text{total}}{K},
\]
que estima a penalidade mínima necessária para que o custo de violar a restrição 
de atribuição única supere o ganho de qualquer atribuição inválida.

Foram totalizadas 504 execuções com duas instâncias. Na instância regular,
usou-se $\mathbf{w}=[0.15, 0.35, 0.25, 0.25]$ e $P_\text{safe}=0.35$.
Na instância com processo dominante, usou-se
$\mathbf{w}=[0.70, 0.10, 0.10, 0.10]$ e $P_\text{safe}=0.70$.

\subsection{Viabilidade sob orçamento reduzido}

Primeiro foi fixado $p=1$, \texttt{top\_k}=1, e o número de passos do 
otimizador entre 1 e 50. Os resultados agregam duas instâncias, duas penalidades 
e cinco sementes aleatórias para inicilização de $\gamma$ e $\beta$, com 20 execuções por mixer em cada orçamento.

\begin{table}[H]
\centering
\begin{tabular}{|c|c|c|c|c|}
\hline
\textbf{Passos} & \textbf{X viável} & \textbf{XY viável} & 
\textbf{X ótimo} & \textbf{XY ótimo} \\
\hline
1  & 6/20  & 20/20 & 6/20  & 0/20 \\
3  & 12/20 & 20/20 & 12/20 & 0/20 \\
5  & 12/20 & 20/20 & 12/20 & 0/20 \\
10 & 20/20 & 20/20 & 20/20 & 0/20 \\
20 & 20/20 & 20/20 & 19/20 & 0/20 \\
50 & 20/20 & 20/20 & 17/20 & 0/20 \\
\hline
\end{tabular}
\caption{Viabilidade e otimalidade com \texttt{top\_k}=1 e $p=1$}
\label{tab:resultados-mx3}
\end{table}

O mixer XY manteve viabilidade em 120/120 execuções e atribuiu sempre 100\% da 
massa de probabilidade ao subespaço de atribuições únicas válidas. O mixer X falhou a restrição 
de atribuição única em 30/120 execuções, com todas as falhas abaixo de 10 passos.

Contudo, o padrão de otimalidade inverteu-se: com \texttt{top\_k}=1, o mixer X 
produziu 86/120 soluções ótimas, enquanto o mixer XY produziu 0/120. A razão é 
estrutural: o XY parte de um estado inicial onde todos os processos estão 
atribuídos ao núcleo 0, e com circuitos pouco profundos a distribuição permanece 
concentrada em atribuições viáveis mas subótimas. O estado ótimo global foi 
classificado em primeiro lugar pelo mixer X nas execuções, mas 
apareceu nas posições 6--10 para o mixer XY na instância de referência.

\subsection{Comparação com Valores Práticos de \texttt{top\_k}}

Para complementar o experimento anterior em que foi usado \texttt{top\_k}=1, foi realizado um segundo
conjunto de experimentos com $p \in \{2, 3\}$, 100 passos, $P_\text{safe}$ e
\texttt{top\_k} $\in \{2, 3, 5, 10\}$. Todas as 160 execuções foram viáveis.

\begin{table}[H]
\centering
\begin{tabular}{|c|c|c|c|c|}
\hline
\textbf{Instância} & \textbf{$p$} & \textbf{\texttt{top\_k}} & 
\textbf{X ótimo} & \textbf{XY ótimo} \\
\hline
Regular   & 2 & 2  & 5/5 & 0/5 \\
Regular   & 2 & 3  & 5/5 & 0/5 \\
Regular   & 2 & 5  & 5/5 & 0/5 \\
Regular   & 2 & 10 & 5/5 & 5/5 \\
Regular   & 3 & 2  & 5/5 & 1/5 \\
Regular   & 3 & 3  & 5/5 & 2/5 \\
Regular   & 3 & 5  & 5/5 & 2/5 \\
Regular   & 3 & 10 & 5/5 & 5/5 \\
Dominante & 2 & 2  & 5/5 & 2/5 \\
Dominante & 2 & 3  & 5/5 & 2/5 \\
Dominante & 2 & 5  & 5/5 & 5/5 \\
Dominante & 2 & 10 & 5/5 & 5/5 \\
Dominante & 3 & 2  & 5/5 & 2/5 \\
Dominante & 3 & 3  & 5/5 & 2/5 \\
Dominante & 3 & 5  & 5/5 & 5/5 \\
Dominante & 3 & 10 & 5/5 & 5/5 \\
\hline
\end{tabular}
\caption{Otimalidade por mixer, profundidade e \texttt{top\_k} com 100 passos e $P_\text{safe}$.}
\label{tab:resultados-mx5}
\end{table}

O mixer X atingiu o ótimo em 80/80 execuções. O mixer XY atingiu 43/80, com
desempenho a melhorar consistentemente com \texttt{top\_k}: ambos convergiram para
100\% de otimalidade em \texttt{top\_k}=10. O tempo médio foi 2787.0\,ms para X e
2200.7\,ms para XY, uma redução de 21\%.

Este resultado clarifica o compromisso observado anteriormente: a dificuldade do XY
em optimalidade não é absoluta, mas depende do número de candidatos disponíveis. Com
\texttt{top\_k}=10, a distribuição do XY contém candidatos suficientes para capturar
o estado ótimo; com valores menores, a concentração de massa em atribuições viáveis
mas subótimas limita a recuperação.

\subsection{Síntese}

Os experimentos revelam um compromisso claro entre os dois mixers. O XY é
estruturalmente superior em viabilidade: manteve 100\% de massa no subespaço de
atribuições únicas válidas em todas as execuções e nunca produziu uma solução
inválida, independentemente da penalidade ou do orçamento do otimizador. Foi também
16--22\% mais rápido em todos os regimes testados.

Em optimalidade, o comportamento depende do número de candidatos disponíveis. Com
\texttt{top\_k}=1, o mixer X superou claramente o XY --- 86/120 ótimos contra 0/120 ---
porque o estado ótimo global tende a aparecer em posições mais elevadas na sua
distribuição. Com valores práticos de \texttt{top\_k}, essa vantagem diminui: a
partir de \texttt{top\_k}=10, ambos os mixers atingiram 100\% de otimalidade.

A escolha do mixer deve portanto considerar o regime de operação. O XY é preferível
quando a viabilidade é prioritária ou quando \texttt{top\_k} é suficientemente grande;
o mixer X é mais robusto em optimalidade quando o orçamento de candidatos é muito
restrito, desde que o otimizador tenha passos suficientes para garantir viabilidade.

\section{Limites da Pipeline}
\label{sec:fase4}

As fases anteriores caracterizaram o comportamento da pipeline em condições favoráveis:
workloads uniformes e instâncias pequenas. Nessas condições, a
pipeline atingiu o ótimo na quase totalidade das execuções. Para identificar onde e como
a qualidade se degrada, foram construídas instâncias deliberadamente difíceis e
alargado o espaço de parâmetros investigado. 

\subsection{Instâncias com Processo Dominante}

A primeira estrutura consiste em seis processos onde um tem peso muito
superior aos restantes. Nestes casos, a atribuição ótima é estruturalmente forçada: o
processo dominante deve ir para um núcleo e os cinco leves para o outro; qualquer outra
distribuição é fortemente penalizada. Foram testados pesos dominantes em
$\{0.60, 0.70, 0.80, 0.90, 0.95\}$, com os restantes uniformes, e combinações
$p \in \{1, 2\}$ e $\mathtt{top\_k} \in \{3, 10\}$.

\begin{table}[H]
\centering
\begin{tabular}{|c|c|c|c|c|}
\hline
\textbf{$p$} & \textbf{\texttt{top\_k}} & \textbf{Ótimos} & \textbf{Gap médio} & \textbf{Tempo médio (ms)} \\
\hline
1 & 3  & 24/25 & 0,006755\% & 2003,0 \\
1 & 10 & 25/25 & 0          & 1999,3 \\
2 & 3  & 9/25  & 0,082395\% & 3004,6 \\
2 & 10 & 25/25 & 0          & 3000,1 \\
\hline
\end{tabular}
\caption{Resultados com processo dominante ($N=6$, $K=2$).}
\label{tab:resultados-e1a}
\end{table}

O resultado mais relevante é que aumentar $p$ de 1 para 2 degradou fortemente a
solução quando $\mathtt{top\_k}=3$: de 24/25 para 9/25 ótimos. 

Com
$\mathtt{top\_k}=10$, a profundidade deixou de fazer diferença e ambas as configurações
atingiram o ótimo em todas as repetições. 

Este padrão revela que maior profundidade não
implica necessariamente melhor qualidade --- em circuitos mais profundos, a distribuição
de probabilidade pode dispersar-se por mais estados, exigindo um conjunto de candidatos
maior para capturar o ótimo.

\subsection{Instâncias com Pesos Quase Iguais}

A segunda estrutura adversarial usa quatro processos com pesos próximos entre si, onde
o desequilíbrio máximo entre o processo mais leve e o mais pesado varia entre 0,01 e
0,20. Estas instâncias são difíceis porque o ótimo global tem energia muito semelhante
a soluções subótimas, tornando a distinção difícil para o otimizador.

\begin{table}[H]
\centering
\begin{tabular}{|c|c|c|c|c|}
\hline
\textbf{$p$} & \textbf{\texttt{top\_k}} & \textbf{Ótimos} & \textbf{Gap médio} & \textbf{Tempo médio (ms)} \\
\hline
1 & 3  & 15/20 & 0,003213\% & 1569,5 \\
1 & 10 & 20/20 & 0          & 1566,0 \\
2 & 3  & 8/20  & 0,005712\% & 2167,5 \\
2 & 10 & 20/20 & 0          & 2167,1 \\
\hline
\end{tabular}
\caption{Resultados com pesos quase iguais ($N=4$, $K=2$).}
\label{tab:resultados-e1b}
\end{table}

O padrão é consistente com o anterior: $\mathtt{top\_k}=10$ eliminou todas as falhas,
enquanto $\mathtt{top\_k}=3$ deixou uma fração de runs subótimos. Os gaps são pequenos
em valor absoluto, mas não equivalem a optimalidade exata. O simulated annealing atingiu
o ótimo em 80/80 execuções neste conjunto.

\subsection{Instâncias com Estrutura CPU/I\slash O}

A terceira estrutura adversarial foi concebida para testar a decomposição adaptativa. Foram criados seis processos divididos em dois grupos
distintos: três CPU-bound (pesos altos, I/O wait baixo) e três I/O-bound (pesos baixos,
I/O wait alto). 

O parâmetro $\mathtt{io\_alpha}$ controla o peso relativo do perfil de
I/O na afinidade da decomposição.

\begin{table}[H]
\centering
\resizebox{\textwidth}{!}{%
\begin{tabular}{|c|c|c|c|c|c|}
\hline
\textbf{\texttt{io\_alpha}} & \textbf{Ótimos} & \textbf{Coupling intra} & \textbf{Coupling inter} & \textbf{Load imbalance} & \textbf{Tempo (ms)} \\
\hline
0,1 & 20/20 & 0,977 & 0,295 & 0,561 & 716,3 \\
0,5 & 20/20 & 0,984 & 0,294 & 0,647 & 715,5 \\
0,9 & 20/20 & 0,987 & 0,293 & 0,732 & 714,0 \\
\hline
\end{tabular}
}
\caption{Resultados do clustering CPU/I/O para diferentes valores de \texttt{io\_alpha} ($N=6$, $K=2$).}
\label{tab:resultados-e1c}
\end{table}

Em todas as configurações, a decomposição separou corretamente os três processos
CPU-bound dos três I/O-bound As 60 execuções atingiram o ótimo.

No entanto, aumentar $\mathtt{io\_alpha}$ aumentou o desequilíbrio de carga: ao
privilegiar a afinidade de I/O na formação dos clusters, os bundles ficam com pesos
totais mais desiguais. Este resultado mostra que afinidade da decomposição e qualidade
de scheduling são objetivos potencialmente conflituosos --- uma decomposição mais pura em
termos de perfil de I/O não implica necessariamente uma atribuição mais equilibrada.
